\section{Empirical Results}
\label{sec:empirical}

\subsection{Responsiveness and Bias Comparison}

Table~\ref{tab:sv-summary} reports the seats-votes curve summary statistics for
\textsc{bisect} and enacted 2022 congressional maps across four states using
VEST 2020 presidential returns.

\begin{table}[ht]
\centering
\caption{Seats-Votes Curve Statistics: Bisect vs.\ Enacted 2022 Congressional Maps}
\label{tab:sv-summary}
\begin{threeparttable}
\begin{tabular}{@{}llrrrr@{}}
\toprule
State ($k$) & Plan & $S(0.50)$ & Bias & $R$ & $S(0.50) \times k$ \\
\midrule
NC ($k=14$) & Bisect$^\dagger$ & $0.500$ & $0.000$ & $2.0$ & $7.0$ \\
NC ($k=14$) & Enacted 2022    & $0.429$ & $-0.071$ & $1.3$ & $6.0$ \\
\addlinespace
WI ($k=8$)  & Bisect$^\dagger$ & $0.500$ & $0.000$ & $1.9$ & $4.0$ \\
WI ($k=8$)  & Enacted 2022    & $0.375$ & $-0.125$ & $1.4$ & $3.0$ \\
\addlinespace
TX ($k=38$) & Bisect$^\dagger$ & $0.500$ & $0.000$ & $2.1$ & $19.0$ \\
TX ($k=38$) & Enacted 2022    & $0.474$ & $-0.026$ & $1.2$ & $18.0$ \\
\addlinespace
FL ($k=28$) & Bisect$^\dagger$ & $0.500$ & $0.000$ & $2.0$ & $14.0$ \\
FL ($k=28$) & Enacted 2022    & $0.464$ & $-0.036$ & $1.3$ & $13.0$ \\
\bottomrule
\end{tabular}
\begin{tablenotes}
\small
\item $^\dagger$ Single-run \textsc{bisect} result. Negative Bias indicates Republican advantage.
$R$ = responsiveness ($dS/dv$ at $v=0.50$). $S(0.50) \times k$ = expected Democratic seats
at 50\% statewide vote share.
\end{tablenotes}
\end{threeparttable}
\end{table}

The NC result is particularly striking: at a 50-50 statewide partisan vote split,
\textsc{bisect} produces 7 Democratic seats and 7 Republican seats; the enacted map
produces 6 Democratic seats and 8 Republican seats. The one-seat difference is the
Partisan Bias, and it is directly visible in the seats-votes curve without any
computation beyond reading $S(0.50)$ from the graph.

Responsiveness differences are equally dramatic: NC bisect $R = 2.0$ vs.\ enacted
$R = 1.3$---a 35\% reduction in electoral accountability. In practical terms, under
the enacted map, a 5-point Democratic surge produces approximately $5 \times 1.3/100
\times 14 \approx 0.9$ additional Democratic seats. Under bisect, the same surge
produces $5 \times 2.0/100 \times 14 \approx 1.4$ additional Democratic seats.
The difference (0.5 seats) represents the incumbent protection embedded in
safe-seat structure of the enacted map.

\subsection{NC Seats-Votes Curve Data}

Table~\ref{tab:nc-sv-data} reports the full NC seats-votes curve at selected
vote-share levels, illustrating the asymmetry between bisect and enacted curves
across the full plausible vote range.

\begin{table}[ht]
\centering
\caption{NC ($k=14$) Seats-Votes Curve: Selected Vote Levels}
\label{tab:nc-sv-data}
\begin{tabular}{@{}lrr@{}}
\toprule
$v$ (D vote share) & $S(v)$ Bisect$^\dagger$ & $S(v)$ Enacted \\
\midrule
0.35 & $0.000$ & $0.000$ \\
0.40 & $0.071$ & $0.000$ \\
0.45 & $0.214$ & $0.143$ \\
0.48 & $0.357$ & $0.286$ \\
0.50 & $0.500$ & $0.429$ \\
0.52 & $0.571$ & $0.500$ \\
0.55 & $0.714$ & $0.571$ \\
0.60 & $0.857$ & $0.643$ \\
0.65 & $0.929$ & $0.786$ \\
0.70 & $1.000$ & $0.929$ \\
\bottomrule
\end{tabular}
\end{table}

The table demonstrates the full seats-votes asymmetry: at every vote level above
40\%, \textsc{bisect} produces equal or higher Democratic seat share than enacted.
The gap is largest near $v = 0.50$--$0.55$, the politically relevant range.

\subsection{Recovering EG from the Curve}

Using the NC data in Table~\ref{tab:nc-sv-data}, EG can be verified as the area
between $S(v)$ and $S(v)=v$ under the uniform distribution on $[0,1]$. Numerically
integrating the enacted curve: the area above the diagonal (where enacted $S(v) > v$,
i.e., for large $v$) minus the area below the diagonal (where enacted $S(v) < v$)
yields approximately $+0.09$, consistent with the EG reported in L.1. The sufficiency
of $S(v)$ for EG is thus confirmed empirically.
